\begin{table}[H]
\small
\setlength{\tabcolsep}{4pt}

\caption{\label{tab:ea20-policy-method-sensitivity}Annual effects of household price policies by counterfactual construction method in the euro area, 2021--2023}
\centering
\begin{tabular}[t]{lrrrrrr}
\toprule
\multicolumn{1}{c}{} & \multicolumn{3}{c}{\shortstack{Effect of price policies on\\average inflation}} & \multicolumn{3}{c}{\shortstack{Effect of price policies on\\Q1--Q5 inflation inequality}}\\
\cmidrule(lr){2-4} \cmidrule(lr){5-7}
Construction method & 2021 & 2022 & 2023 & 2021 & 2022 & 2023\\
\midrule
All price measures & -0.1 & -1.5 & -0.6 & -0.1 & -0.6 & -0.5\\
Directly reconstructed measures only & 0.0 & -0.8 & 0.0 & 0.0 & -0.3 & -0.1\\
\bottomrule
\end{tabular}
\begin{flushleft}
\footnotesize{Notes. Observed minus counterfactual annual inflation or Q1--Q5 gap, in percentage points; negative values indicate a reduction. Directly reconstructed measures use statutory rates, unit rebates, or published reference-price paths without a fitted intensity. All price measures additionally include calibrated paths and conditional German and Slovenian bill or tariff benchmarks. Sources: Eurostat, national sources, authors' calculations.}
\end{flushleft}
\end{table}
